% FENG 498 Defense -- Anticipated jury questions for the Methodology & Simulation part.
% Deniz Ege Memetoglu.  Compile: pdflatex qa_deniz_ege.tex
% Assumes the jury has the written report and the presentation in front of them; each
% answer ends with a pointer (->) to the relevant report section / slide / figure.
\documentclass[11pt]{article}
\usepackage[utf8]{inputenc}
\usepackage[T1]{fontenc}
\usepackage{lmodern}
\usepackage[margin=2.1cm]{geometry}
\usepackage{amsmath}
\usepackage{xcolor}
\usepackage{enumitem}
\usepackage{parskip}
\usepackage{titlesec}

\definecolor{navy}{HTML}{0F2747}
\definecolor{green}{HTML}{218A3A}
\definecolor{grey}{HTML}{5B6672}

\titleformat{\section}{\color{navy}\large\bfseries}{}{0pt}{}[{\color{navy}\titlerule[1pt]}]
\titlespacing{\section}{0pt}{12pt}{6pt}

\newcommand{\q}[1]{\par\medskip\noindent{\color{navy}\bfseries Q.\ #1}\par\smallskip}
\newcommand{\pt}[1]{\par{\color{grey}\small$\hookrightarrow$\ #1}\par}

\begin{document}
\thispagestyle{empty}

\begin{center}
  {\color{navy}\LARGE\bfseries Anticipated Jury Questions}\\[2pt]
  {\large Methodology \& Simulation (Section 3)}\\[4pt]
  {\color{grey} Deniz Ege Memeto\u{g}lu \,\textbullet\, FENG 498 \,\textbullet\, assumes the jury has the report and the presentation in front of them}
\end{center}
\medskip
{\color{navy}\rule{\linewidth}{1.5pt}}

% ───────────────────────────────────────────────
\section{Model design}

\q{Is this a trace replay of the logs, or a real model?}
Distribution-driven, not a replay. I fit distributions (inter-batch gap, batch size, kit
composition, service times) from the four-month ZWM92 dispatch log, 167{,}784 rows and
40{,}804 reconstructed kit orders, and the model samples from them. It does not play
history back. That is deliberate: because demand is generated and not fixed, I can move
materials to new locations and re-run the same warehouse under a different policy, which
a replay could never do. It is built in Python with SimPy.
\pt{Report: Methodology, Simulation Model (Modelling Approach, Order Generation). Slide: the methodology flowchart and the distribution-driven logic slide.}

\q{Why discrete-event simulation, and why SimPy?}
The bottleneck is queueing for shared resources: operators waiting on a seven-truck
reach-truck fleet, plus the cascade where one order holds its slot while it waits and
delays the next. Static or analytic models assume average flow and miss that contention.
SimPy is a mature, transparent discrete-event library; every resource, operators, reach
trucks, Kardex units and pallet-position locks, is an explicit, auditable object.
\pt{Report: Methodology, Simulation Model. Slide: the SimPy resource-model figure.}

% ───────────────────────────────────────────────
\section{Inputs and calibration}

\q{Where do the pick and handling times come from?}
From a video time-motion study of the F400 kitting line: 2{,}319 timed micro-events over
296.6 observed minutes. Service times are lognormal, with the spread set by the measured
coefficient of variation through \(\sigma=\sqrt{\ln(1+\mathrm{CV}^2)}\). The means are
operator pick 0.113~min, reach-truck pick-and-place 0.110, manual-band penalty 0.102, and
Kardex pick 0.113; the sigmas are 1.245, 1.047 and 1.279.
\pt{Report: Methodology, Data Collection (Time and Motion Study) and Simulation Model (Travel, Picking and Service Times). Slide: the sigma formula on the logic slide.}

\q{How did you validate the Kardex timing? There was no Kardex video.}
Honest answer: there was no Kardex-specific footage, so the Kardex pick mean and CV mirror
the F400 operator pick, with a 0.4-minute carousel turn as a book value. We disclose this
in the assumption register. It cannot bias the comparison, because Kardex routing is
identical across all six policies, so any timing error cancels out of the ranking.
\pt{Report: Appendix, Model Assumption Register.}

\q{F400 timing is extrapolated to nine families. What if F400 is unrepresentative?}
F400 is about 29\% of the dispatch rows, and its timing is applied to the other eight
families. We bound that risk with a one-at-a-time sensitivity sweep of plus or minus
twenty percent on every timing input. Line-aware stays the best policy under every swing;
lead time moves proportionally, but the ranking never flips. Per-line video studies are
listed as future work.
\pt{Report: Results, Sensitivity Analysis (tornado figure).}

\q{Why model batch arrivals, and does the rate match reality?}
Because the real log arrives in bursts: 3{,}885 same-timestamp batches, mean 4.68 kits per
batch, with an inter-batch gap mean of 5.36 minutes and a CV of 2.04, so it is empirical,
not Poisson. Across the full twenty-replication set the model produces about 400 orders a
day, within 1.1\% of the real 396-per-day average.
\pt{Report: Methodology, Simulation Model (Order Generation); Results, Verification and Validation Results (daily-volume calibration).}

% ───────────────────────────────────────────────
\section{Experiment design}

\q{Why twenty replications, and why the same seeds across policies?}
Each policy runs twenty replications; replication \(i\) uses seed \(42+i\), with separate
random streams for arrivals and for service times. Common random numbers mean every policy
faces exactly the same demand, so the differences are the policy, not luck. The 95\%
confidence intervals, using the t-distribution with 19 degrees of freedom, are narrow and
do not overlap between Line-aware and the baselines, which is why twenty runs are enough.
\pt{Report: Methodology, Experimental Design and Statistical Analysis (replication analysis). Slide: the statistics slide.}

\q{Why six policies, not only your proposed one?}
To place the proposal against a fair spread: the real SAP placement as the operational
baseline, a simple capacity heuristic as a control, two literature-based ABC variants, a
pure travel-distance optimizer, and our Line-aware policy. Without those references,
``better'' would mean nothing. The policy design itself is Nehir's section, but the
experimental frame is shared.
\pt{Report: Methodology, Storage Assignment Policies.}

% ───────────────────────────────────────────────
\section{Verification and validation}

\q{Your chi-square is highly significant. Doesn't that mean the model is wrong?}
No, and this is the key point. The pooled chi-square with 7 degrees of freedom is 535,
p below 0.001, but Cramer's V is only 0.22, a small-to-moderate effect. It is significant
precisely because the model uses its own routing logic, not a bin-for-bin replay; an exact
match would only happen if we replayed history, which would defeat the purpose. We read the
effect size, and we treat it as an internal-consistency check.
\pt{Report: Results, Verification and Validation Results. Slide: the validation slide with the chi-square and Cramer's V formula.}

\q{You fitted and validated on the same ZWM92 data. No holdout?}
Correct, and we disclose it. SAP never logged waiting times or reach-truck calls, so there
was no independent observation period to hold out. The independent grounding is the F400
video for timing, the daily-volume calibration to within about one percent, and a 3D route
audit in which no agent crosses a rack body. We label these as internal-consistency tests,
not out-of-sample validation.
\pt{Report: Methodology, Verification and Validation Approach; Results, V\&V Results.}

% ───────────────────────────────────────────────
\section{Limitations and scope}

\q{Aisle congestion is not modelled. Doesn't that invalidate the results?}
It makes the absolute lead times slightly optimistic under heavy load, but all six policies
share the assumption symmetrically, so the comparison stays fair. If anything, adding
congestion would widen Line-aware's lead, because the baselines saturate the reach-truck
fleet far more, up to 44\% utilization for the travel-distance policy.
\pt{Report: Appendix, Model Assumption Register; Conclusions, Limitations.}

\q{How real is the ``Actual SAP'' baseline?}
It is an honestly-labelled SAP-plus-heuristic hybrid. 750 materials sit at their exact
decoded SAP rack-bay-position, 2{,}872 are Kardex-routed and policy-invariant, and the rest
fall to a frequency heuristic. 386 multi-bin materials are not split into forward and
reserve. We never claim a pure replay; it is the best reconstruction the decoded SAP data
allows, and it is our validation ground truth.
\pt{Report: Methodology, Data Collection; Conclusions, Limitations.}

\q{Throughput barely changes and RT utilization is only 1.3\%. Is that believable?}
Yes. The system is demand-limited: every policy sees the same arrival stream, so throughput
is about 398 a day for all of them, with an ANOVA F near 0.01. The win is in lead and
waiting time, not output. The 1.3\% reach-truck figure is picking workload only;
replenishment is excluded for all six policies equally, so the comparison is fair, and we
flag that the live total would be higher.
\pt{Report: Results, Policy Comparison and Discussion; Conclusions, Limitations (replenishment caveat).}

\q{What about the milk-run trains?}
The model represents seven milk-run trains as an available delivery mode, but in the
reported runs the operator carries the finished kit to the line, so lead time is attributed
to operators and reach trucks only. Enabling milk-run would offload some operator walking;
it does not touch the reach-truck queue that drives the result.
\pt{Report: Appendix, Model Assumption Register.}

% ───────────────────────────────────────────────
\section{Layout and reproducibility}

\q{Is the warehouse layout real or invented?}
Reconstructed from the real AutoCAD drawing and eleven per-rack pallet PDFs, cross-checked
by a blind CAD extraction. Eleven rack rows, A to J and U; 3{,}101 modelled pallet positions
reconcile with the 3{,}203 in the drawing, the 102-gap being narrower feeder bays plus a
cart-parking slot.
\pt{Report: Methodology, System Description and Geometry; the top-down layout figure.}

\q{Can someone reproduce your results?}
Yes. The code is fixed-seed and deterministic given a seed; six policies times twenty
replications is 120 runs, exported as a tidy table for Minitab. The software versions and
reproducibility notes are in the appendix.
\pt{Report: Appendix, Software and Reproducibility.}

\vfill
{\color{grey}\small One-line fallbacks if you blank: ``It is distribution-driven, so I can
re-slot and re-run.'' \,\textbullet\, ``Chi-square is significant by design; read Cramer's
V, 0.22.'' \,\textbullet\, ``1.3\% RT is picking only, and it is fair across all policies.''
\,\textbullet\, ``Sensitivity at $\pm$20\% keeps Line-aware first.''}

\end{document}
